\chapter{Regressão com Descontinuidade}
\label{cap:rdd}
% ── índice ──────────────────────────────────────────────────────
\index{regressão por descontinuidade|textbf}
\index{RDD|see{regressão por descontinuidade}}
\index{rd@\texttt{rd()}|textbf}
\index{limiar de corte}
\index{running variable}

% ─────────────────────────────────────────────────────────────────────────────
\section{O modelo}

A Regressão com Descontinuidade (RDD) identifica efeitos causais quando o
tratamento é atribuído mecanicamente com base em um limiar de uma variável
contínua observável — a \textbf{variável de corrida} (\emph{running variable})
$x$:

\[
  D_i = \mathbf{1}[x_i \geq c]
\]

O estimador explora o fato de que indivíduos imediatamente abaixo e acima do
limiar $c$ são comparáveis em tudo, exceto no tratamento recebido. O efeito
identificado é o \textbf{LATE no limiar} (\emph{Local Average Treatment
Effect at the cutoff}):

\[
  \tau_{\mathrm{RDD}}
  = \lim_{x \downarrow c} E[y \mid x]
  - \lim_{x \uparrow c} E[y \mid x]
\]

\begin{nota}
  O RDD \emph{Sharp} (acima) assume que $D$ muda deterministicamente em $c$.
  O RDD \emph{Fuzzy} generaliza para quando $D$ aumenta
  \emph{descontinuamente} em $c$, sem salto de 0 a 1 — requer um primeiro
  estágio e equivale a um IV local.
\end{nota}

% ─────────────────────────────────────────────────────────────────────────────
\section{Identificação}

A hipótese central é a \textbf{continuidade dos resultados potenciais} em
$c$:

\[
  E[Y^{(0)} \mid x] \quad \text{e} \quad E[Y^{(1)} \mid x]
  \quad \text{são contínuas em } x = c
\]

Isso garante que qualquer descontinuidade em $E[y \mid x]$ no ponto $c$ é
causada pelo salto em $D$, e não por outro fator que também muda em $c$.

A hipótese é violada se agentes \emph{manipulam} $x$ para ficarem acima ou
abaixo do limiar — estudantes que revisam respostas para cruzar o corte de
aprovação, por exemplo. O teste de densidades de McCrary (teste de
continuidade da densidade de $x$ em $c$) verifica se há aglomeração
suspeita de observações de um lado do limiar.

% ─────────────────────────────────────────────────────────────────────────────
\section{O estimador}

O estimador padrão é a \textbf{regressão local linear} separada em cada
lado do limiar, dentro de uma janela (\emph{bandwidth}) $h$:

\[
  \hat\tau = \hat\mu_+ - \hat\mu_-
\]

onde $\hat\mu_+$ e $\hat\mu_-$ são os interceptos das regressões locais
lineares no lado direito ($x \in [c, c+h)$) e esquerdo ($x \in (c-h, c)$),
respectivamente, ponderadas pelo kernel triangular:

\[
  K\!\left(\frac{x - c}{h}\right) = \left(1 - \frac{|x-c|}{h}\right)
  \mathbf{1}[|x-c| \leq h]
\]

O \emph{bandwidth} $h$ controla o viés-variância: janelas largas reduzem
variância mas aumentam viés (linearidade local pode ser má aproximação);
janelas estreitas são mais locais mas têm menos observações. O Hayashi
seleciona $h$ automaticamente pelo critério MSE-ótimo (Imbens-Kalyanaraman)
quando \lstinline{bw} não é especificado.

% ─────────────────────────────────────────────────────────────────────────────
\section{Verificação por DGP}

\subsection*{Recuperação exata (sem ruído)}

O DGP usa a variável de corrida $x \sim \mathrm{Uniforme}(-1, 1)$, limiar em
$c = 0$ e efeito verdadeiro $\tau = 2{,}0$. A função de resultados é linear
em $x$ dos dois lados do limiar, sem ruído, de modo que a RDD deve recuperar
$\tau$ exatamente.

\[
  y_i = 1{,}0 + 0{,}5\,x_i + 2{,}0\,D_i, \qquad D_i = \mathbf{1}[x_i \geq 0]
\]

\begin{lstlisting}[caption={DGP RDD sem ruído — recuperação exata}]
seed(42)
input df
id
1
end
for i in 1..=10 { df = append(df, df) }
df |> filter(_n <= 1000)
generate df x = runiform() * 2 - 1
generate df d = x >= 0
generate df y = 1.0 + 0.5*x + 2.0*d
rd(y ~ x, 0.0, df)
\end{lstlisting}

\begin{lstlisting}[language={}, basicstyle=\ttfamily\small,
  frame=none, numbers=none, backgroundcolor=\color{codbg},
  xleftmargin=0pt, xrightmargin=0pt, breaklines=false]
══════════════════════════════════════════════════════════════════════
 Regressão Descontínua  —  Sharp  —  Local Linear (p=1)
══════════════════════════════════════════════════════════════════════
 Outcome: y                   Running var: x
 Cutoff: 0.0000   Bandwidth: 0.0571   Kernel: Triangular
 Obs total: 72   N esquerda: 37   N direita: 35
──────────────────────────────────────────────────────────────────────
 Efeito de Tratamento (τ̂):
       2.0000   SE     0.0000   z 2.443e15   P>|z|   0.0000  ***
 IC 95%: [2.0000, 2.0000]
══════════════════════════════════════════════════════════════════════
 *** p<0.01  ** p<0.05  * p<0.10
\end{lstlisting}

O estimador recupera $\hat\tau = 2{,}0000$ exatamente. O bandwidth de
$0{,}0571$ foi selecionado automaticamente — das 1000 observações, 72 estão
dentro da janela em torno do cutoff.

\subsection*{Com ruído idiossincrático}

Adicionando $\varepsilon$ idiossincrático, o LATE permanece identificado:

\begin{lstlisting}[caption={DGP RDD com ruído}]
generate df e = rnormal()
generate df y = 1.0 + 0.5*x + 2.0*d + e
rd(y ~ x, 0.0, df)
\end{lstlisting}

\begin{lstlisting}[language={}, basicstyle=\ttfamily\small,
  frame=none, numbers=none, backgroundcolor=\color{codbg},
  xleftmargin=0pt, xrightmargin=0pt, breaklines=false]
══════════════════════════════════════════════════════════════════════
 Regressão Descontínua  —  Sharp  —  Local Linear (p=1)
══════════════════════════════════════════════════════════════════════
 Outcome: y                   Running var: x
 Cutoff: 0.0000   Bandwidth: 0.0741   Kernel: Triangular
 Obs total: 94   N esquerda: 43   N direita: 51
──────────────────────────────────────────────────────────────────────
 Efeito de Tratamento (τ̂):
       2.0116   SE     0.1191   z   16.893   P>|z|   0.0000  ***
 IC 95%: [1.7782, 2.2449]
══════════════════════════════════════════════════════════════════════
 *** p<0.01  ** p<0.05  * p<0.10
\end{lstlisting}

Com ruído, $\hat\tau = 2{,}0116$ — o desvio de $2{,}0$ é variância amostral.
O IC de 95\% $[1{,}78;\; 2{,}24]$ cobre o verdadeiro $\tau = 2{,}0$.
O bandwidth foi $0{,}0741$: com mais variância, o critério MSE-ótimo
prefere janela mais estreita para controlar o viés.

% ─────────────────────────────────────────────────────────────────────────────
\section{Sintaxe}

\begin{lstlisting}
rd(y ~ x, cutoff, df)
rd(y ~ x, cutoff, df, bw=0.3)
rd(y ~ x, cutoff, df, bw=0.3, poly=2)
\end{lstlisting}

\begin{center}
\begin{tabular}{lll}
\toprule
\textbf{Parâmetro} & \textbf{Padrão} & \textbf{Descrição} \\
\midrule
\texttt{cutoff}  & obrigatório & valor do limiar em $x$ \\
\texttt{bw}      & automático (IK) & largura da janela $h$ \\
\texttt{poly}    & \texttt{1} & grau do polinômio local \\
\texttt{kernel}  & \texttt{triangular} & \texttt{uniform}, \texttt{epanechnikov} \\
\bottomrule
\end{tabular}
\end{center}

\begin{lstlisting}[caption={Fluxo típico — nota de corte e bolsa}]
load "escola.csv" as df

// pontuação centralizada no limiar 60
generate df xc = nota - 60

let m = rd(aprovado ~ xc, 0.0, df)
display m
\end{lstlisting}

\section{Capturando o resultado}

\begin{lstlisting}
let m = rd(y ~ x, 0.0, df)
display m
\end{lstlisting}

\begin{nota}
  Em aplicações reais, valide o RDD com testes de falsificação:
  (1)~estime o modelo usando \emph{covariáveis pré-tratamento} como
  desfecho — não deve haver descontinuidade em $c$;
  (2)~teste descontinuidades em \emph{limiares placebo} ($c \pm \delta$)
  — não devem ser significativas; (3)~verifique a densidade de $x$ em
  $c$ (McCrary) para detectar manipulação da variável de corrida.
\end{nota}
